% !TEX program = xelatex
%
% ENGINE. This deck sets EB Garamond and Hiragino Sans through fontspec, which only XeTeX
% and LuaTeX provide. pdflatex cannot build it at all: it stops at fontspec.sty with
% "requires either XeTeX or LuaTeX" and writes no PDF. The line above tells an editor or
% latexmk which engine to use; slides/.latexmkrc says the same to a bare `latexmk` run.
% Building by hand is `bash slides/build.sh`, which calls xelatex twice.
%
% PFE defense, 20 minutes.
%
% ORDER. The deck runs in the manuscript's own order and nowhere else: introduction, methods in
% dependency order, numerical experiments, the polyester sheet. A talk that regroups the paper by
% theme reads as a different argument from the one being defended, so each part below names the
% section of paper/sn-article.tex it follows. It ends on the measured sheet: a bulleted recap of
% what was just shown spends a minute of twenty saying it again.
%
% THE CLAIM RULE. \punch is for a slide that asserts a result, and most slides do not. A slide
% showing what a thing IS takes a \framesubtitle and no rule. Put one on every frame and every
% frame has to have a thesis, which reads as a deck arguing with the room rather than a person
% explaining their work, and it costs the plate 17 to 41 pt of height per frame besides.
%
% PANEL LETTERS AND FRAMES round a plate are print furniture: a letter is there so running prose
% can write "see (b)", and a rule groups columns because a page has no one to say which are the
% baselines. Here there is a speaker. slides/make_slide_figures.sh builds the plates that carry
% them under NFF_SLIDE_FIGURES=1, read by figstyle.panel_letter, plus --no-boxes where it
% applies. The paper's own copies in paper/Images are untouched and still lettered.
%
% Every number on a slide appears in the LIVE prose of paper/sn-article.tex, not in a comment.
% A number the manuscript has retired is deleted here rather than updated: a figure that survives
% only as commented-out bookkeeping has no source. Nothing is rounded up, nothing is carried over
% from an earlier draft. Credit lines are the exception: a volume, an article number and a year
% are provenance for a borrowed picture, not a claim this work is making.
%
% The vocabulary follows the manuscript's own retirements: sheet material not stock, internal
% force distribution not force network, training domain not trust region or fitted range, hinge
% design vector not descriptor, boundary position not boundary slider, panel not tile, path not
% trajectory, separation not departure. No em dash in rendered text.
%
%   bash slides/build.sh
%
\documentclass[aspectratio=169,11pt]{beamer}
\usetheme{charter}
\usepackage{amsmath}
\usepackage{bm}
% Backup frames sit after the closing slide and must not inflate the count the audience reads
% on the frames that are actually presented: this restarts the frame number at \appendix.
\usepackage{appendixnumberbeamer}

% EB Garamond carries no CJK, and the craft this work is named after is written in it. Hiragino
% is the macOS system face; if the deck is ever built elsewhere, drop the \jp uses rather than
% letting the glyphs fall back silently to tofu.
% ⚠ NOTHING CALLS THIS any more: the only use was the opening slide's Ise-katagami subtitle,
% and that slide was deleted 2026-08-24. Kept so the deck can name the craft again in kanji.
\newfontfamily\jp{Hiragino Sans}[UprightFont={* W3}, BoldFont={* W6}]
\graphicspath{{figures/}{assets/}{assets/kirigami/}{assets/apps/}{assets/photos/}}

% A caption under one panel of a row: what the panel is, and nothing else.
\newcommand{\platelabel}[1]{\\[0.9mm]{\scriptsize\color{chartergrey}#1}}

% The three states, and the two charts, cut out of the merged validation plate so each gets the
% full slide width. The plate is 370.6 x 438.1 bp and the blank bands between its three rows are
% centred at 164.82 and 292.45 bp from the top, measured on a 150 dpi rasterisation of the build
% that slides/make_slide_figures.sh produces: --no-boxes, and NFF_SLIDE_FIGURES=1 for the letters.
% ⚠ These offsets are measured on the LETTERLESS build and do not transfer to the paper's own
%   plate. With letters the bands sat some 8 bp higher, because a letter is drawn above its row.
% Cutting is trim/clip, not a raster crop: the vector content is untouched.
% ── Symbols, colour-keyed to the plates ────────────────────────────────────────────────────
% A symbol set on a slide carries the SAME colour it carries inside the drawings, so the symbol
% read in a sentence and the symbol found in a figure are visibly one thing. The keys are the
% generators' own, not chosen here: r is R_SLIDER = PURPLE and tau is TAU_SLIDER = GREEN in
% plot_pipeline_overview.py, hinges are teal throughout, loads and targets are red.
% ⚠ Do not recolour these to taste. They are the legend for every plate in the deck.
\newcommand{\symcut}[1]{\textcolor{charterpurple}{$#1$}}
\newcommand{\symbnd}[1]{\textcolor{chartergreen}{$#1$}}
\newcommand{\symhng}[1]{\textcolor{charterteal}{$#1$}}
\newcommand{\symtgt}[1]{\textcolor{charterred}{$#1$}}

\newcommand{\physstates}[1][\charterplatew]{%
  \includegraphics[trim=0 145.63bp 0 164.82bp,clip,width=#1]{physical_validation_slide.pdf}}
\newcommand{\physcharts}[1][\charterplatew]{%
  \includegraphics[trim=0 0 0 292.45bp,clip,width=#1]{physical_validation_slide.pdf}}
% The top row of the same plate: the deployed sheet in plan, and the one hinge photographed alone
% against its prediction. Cut from the merged plate rather than staged as the standalone
% hinge_buckle.pdf that plot_hinge_buckle_figure.py still builds: that generator is the pre-merge
% half, and its panel (a) predates the unmirroring, the squaring of the taped target and the
% 500 mm bar. Everything the manuscript publishes for this content is in the row below.
% The cut is the blank band at 164.82 bp, the same one \physstates starts from: 438.1 - 164.82 =
% 273.26. On the lettered plate it could not be, because the letters of the next row are drawn
% above that band and a cut through it handed this row the next row's labels. Without them the
% band is simply the row boundary, and this row gains the 8 bp the dodge used to cost it.
% Unlike the two above this one takes a HEIGHT fraction, as \wideplate does: the row is 2.25
% wide to 1 tall, so at full slide width it is the height that binds, not the width.
\newcommand{\physhinge}[1][0.82]{%
  \charterplateheight
  \begin{center}%
    \includegraphics[trim=0 273.26bp 0 0,clip,%
                     width=\charterplatew,height=#1\charterbodyht,keepaspectratio]%
                    {physical_validation_slide.pdf}%
  \end{center}%
}

% The two rows of the hinge-path plate, cut apart so each gets the full body width. The plate is
% a 2 x 2 grid -- rows are the two designs, columns are hinge activation and the strain
% trajectory -- over a shared bottom strip carrying the key and the activation colour bar. In a
% plateside column the plate printed at 30 per cent of the size it was authored at: that column is
% 207 pt wide against a plate authored 525.6 bp wide. A row is 2.4 wide to 1 tall, so cut apart
% each row is width-bound on the full body instead.
%
% The plate is 525.6 x 430.992 bp and the cut is at 217.42 bp from the top, measured by scanning a
% 150 dpi rasterisation of figures/hinge_path_design.pdf for blank rows.
% ⚠ The cut is NOT the widest blank band, which is at 201.10. Row 1's x-axis LABEL sits BELOW its
%   tick labels, so a cut at 201.10 strands eta_a = a / w_lig at the top of row 2 and leaves row 1
%   without it. 217.42 is the band underneath that label.
% ⚠ Row 1 cannot carry the activation colour bar. It is drawn once, in the bottom strip, and a
%   single rectangular trim cannot hold row 1 and that strip while dropping row 2. Row 1's claim
%   line names the shading in words instead.
% ⚠ The two fractions are NOT equal. The rows are 217.42 and 213.58 bp tall and both are
%   height-bound here, so equal fractions would print row 2 two per cent larger and the click
%   between them would read as a zoom. 0.76 x 213.58 / 217.42 = 0.7466 puts both at one scale.
\newcommand{\hpdrowone}[1][0.76]{%
  \charterplateheight
  \begin{center}%
    \includegraphics[trim=0 213.58bp 0 0,clip,%
                     width=\charterplatew,height=#1\charterbodyht,keepaspectratio]%
                    {hinge_path_design.pdf}%
  \end{center}%
}
\newcommand{\hpdrowtwo}[1][0.7466]{%
  \charterplateheight
  \begin{center}%
    \includegraphics[trim=0 0 0 217.42bp,clip,%
                     width=\charterplatew,height=#1\charterbodyht,keepaspectratio]%
                    {hinge_path_design.pdf}%
  \end{center}%
}

% Card (c) of the manuscript's Figure 1: the deployed sheet, the clamped centre, the applied
% moment, and every hinge coloured by the energy it stores. Both terms of Pi on a real sheet.
% The card BORDER is left out on purpose. Its ink runs at 313.1-314.1 and 426.8-428.7 bp, and the
% crop stops short of both: a cut that keeps three sides of a rounded box leaves it open-ended.
% Measured on a 150 dpi rasterisation of the 437.31 x 424.57 bp plate, not derived from the
% generator's constants: the leaders are drawn OUTSIDE the axes box, so the declared band
% (314.9 to 394.3) shaves the top of "applied moment M". Ink is 337.1-422.9 x 319.5-396.7.
\newcommand{\pipec}[1][0.62]{%
  \includegraphics[trim=334.0bp 314.5bp 12.31bp 27.17bp,clip,%
                   width=\linewidth,height=#1\charterbodyht,keepaspectratio]{pipeline_overview.pdf}}

% Crops of the five-case build. matrix_all.pdf is 432 x 456.04 bp under --no-headers and
% --no-row-labels, and its five case rows and its legend are separated by blank bands. The
% optimised cut pattern is the second drawn column and spans 96.9 to 172.6 bp. All of it measured by scanning a 150 dpi rasterisation
% of THAT file: the label gutter is sized to the longest case name in a build, so a build with
% fewer cases puts its columns somewhere else and these numbers do not transfer.
% ⚠ These are the --no-headers numbers. With the header band the canvas is 437.26 tall and every
%   row boundary below it sits some 60 bp lower, so the two sets are not interchangeable. The
%   row labels move them again: the gutter is sized from them and panel_w takes what it leaves,
%   so without them panel_h goes from 63.8 bp to 79.6 and every boundary moves.
% The row pitch is 79.6 bp, which is panel_h in plot_inverse_design_matrix.py. The boundary
% between rows 1 and 2 has ink crossing it and shows no blank band, so it is taken as the measured
% 164.48 less one pitch.
%
% Cutting with trim/clip rather than re-running the generator five times keeps every case at one
% scale and one type size, which is the whole point of a walk-through: only the row changes.
\newcommand{\mlegend}{\includegraphics[trim=0 0 0 410.0bp,clip,width=\charterplatew]{matrix_all.pdf}}
\newcommand{\mrow}[2]{\includegraphics[trim=0 #2 0 #1,clip,width=\charterplatew]{matrix_all.pdf}}
\newcommand{\mthumb}[2]{%
  \includegraphics[trim=97.5bp #2 260.0bp #1,clip,width=\linewidth]{matrix_all.pdf}}
% A plate wider than the text block has to be shifted left by half its overhang or it runs off
% the right edge. \fullplate does this for a whole plate; these are cut plates placed by hand, so
% they take the same shift through \bled.
\newcommand{\bled}[1]{%
  \par
  \moveleft\charterbleed\vbox{\hsize=\charterplatew\centering #1\par}%
  \par}

% One case already walked through, kept small so the current one is the only large thing.
\newcommand{\mseen}[3]{%
  \begin{minipage}[t]{0.095\textwidth}\centering
    \mthumb{#1}{#2}\\[0.3mm]{\tiny\color{chartergrey}#3}
  \end{minipage}}

% One row of the two-column detail build: the loading diagram and the force transmission, alone.
% matrix_detail.pdf is 432 x 947.44 bp under --columns ad --no-headers --no-row-labels. With two
% columns sharing the whole measure the force panel is 245.7 bp wide against 110 in the
% four-column build, and ws inside the generator is tied to panel_w, so its load-path weights are
% drawn for that width rather than blown up from a narrower one.
% Boundaries measured by scanning a 110 dpi rasterisation: pitch 177.86, four of the five taken
% from the blank band centres and the first extrapolated back one pitch.
% ⚠ These are matrix_detail's numbers. matrix_all is a different canvas and a different pitch.
\newcommand{\mdrow}[3][0.87]{%
  \includegraphics[trim=0 #3 0 #2,clip,%
                   width=\charterplatew,height=#1\charterbodyht,keepaspectratio]{matrix_detail.pdf}}

% What the force plates draw, in their own colours. Hand-set rather than cropped from the plate's
% own key band: that band is 52 bp tall and shrinking it to a readout would take its 6.6 pt type
% under 2 pt.
\newcommand{\forcekey}{%
  \tikz[baseline=-0.55ex]{\draw[charterorange,line width=1.5pt] (0,0) -- (4.5mm,0);}\,tension route
  \quad
  \tikz[baseline=-0.55ex]{\fill[charterteal] (0,-0.9mm) rectangle (3.4mm,1.1mm);}\,compression
  \quad
  \tikz[baseline=-0.55ex]{\fill[chartershear] (0,0.1mm) circle (0.85mm);}\,shear interface}

\title{Multiscale Differentiable Physics for the Inverse Design of Kirigami Sheets}
\author{Julien Kloers}
\charterinstitute{Projet de Fin d'\'Etudes~\textperiodcentered~\'Ecole Nationale des Ponts et Chauss\'ees\\[1.2mm]
Form-Finding Lab, Princeton University\\[1.2mm]
\footnotesize S.~Adriaenssens \quad R.~Pastrana \quad I.~Moreira de Oliveira \quad R.~Oval}

\begin{document}

\chartertitlepage

% No standalone contents frame. Every section divider IS the contents, with the section being
% entered picked out of it, so a separate outline right after the title page shows it twice.

% ═══════════════════════════════════════════════════════════════════════════════════════════
\section{Two scales in a cut sheet}
% ═══════════════════════════════════════════════════════════════════════════════════════════
\sectionslide{}

% ⚠⚠ (a) IS NOT A CUT SHEET. It is 3D-printed PLA units joined by separate polyester shims, so
% its hinge is a COMPONENT, not a remainder. That is the distinction this slide is standing on,
% which is why the title and the claim line below no longer assert that all three are cut, and
% why the claim now names the contrast instead. Added on the author's instruction 2026-08-24.
%
% ⚠⚠ (a) IS ALSO NOT CC BY. Nature Materials is closed access and the arXiv preprint carries the
% arXiv non-exclusive licence, which is not a reuse grant. It is reproduced here under the
% quotation/teaching exception for an academic defense, with attribution, and MUST NOT be carried
% into the manuscript or any public posting without written permission from Springer Nature.
% Panel (d) of its Figure 1 is used WHOLE, header chip excluded: cutting through the source's own
% labels is not acceptable.
%
% Three images at ONE printed height. With only three the HEIGHT CAP binds rather than the width
% sum, so each box is aspect_k x 0.78\charterbodyht (= 147.6 pt) with a little slack, and
% keepaspectratio settles it. geom 1.025, units 1.093, robot 0.992.
% ⚠⚠ THIS SLIDE CARRIES NO CREDIT LINE, AND THAT IS A DECISION ABOUT A LIVE TALK ONLY.
%   A citation set in italics and bold beside three plain grey plate labels did not belong to the
%   same typography as the rest of the frame, and it read as a paragraph rather than as a caption.
%   The three \platelabel lines are the caption; the borrowed-source record is assets/SOURCES.md.
%   ⛔ Two of these three images carry obligations that a spoken defense discharges and a shared
%   file does not. app_robot.png is CC BY 4.0, whose attribution requirement travels with the
%   work wherever it goes. app_units.png is closed access, shown under the teaching and quotation
%   exception, and SOURCES.md already forbids it in any public posting. If this deck is ever
%   handed out, posted or exported rather than presented, the credits go back on before it goes.
%
% The katagami stencil opens the row. It is the one image kept from the deleted opening slide, and
% it is here rather than there because it belongs with the others: a sheet whose remaining material
% is the whole of what holds it together. Its credit moved with it, as SOURCES.md requires.
%
% ⚠ The credits name the objects rather than lettering them. Adding an image renumbers every
% letter after it, and a credit that says "(b)" then attributes the wrong picture.
\begin{frame}{What cut sheets are used for}
  \centering
  \begin{minipage}[t]{0.288\textwidth}\centering
    \includegraphics[width=\linewidth,height=0.78\charterbodyht,keepaspectratio]{geom.jpg}
    \platelabel{kimono dyeing stencil}
  \end{minipage}\hfill
  \begin{minipage}[t]{0.307\textwidth}\centering
    \includegraphics[width=\linewidth,height=0.78\charterbodyht,keepaspectratio]{app_units.png}
    \platelabel{impact absorber}
  \end{minipage}\hfill
  \begin{minipage}[t]{0.279\textwidth}\centering
    \includegraphics[width=\linewidth,height=0.78\charterbodyht,keepaspectratio]{app_robot.png}
    \platelabel{crawling robot}
  \end{minipage}
\end{frame}

\begin{frame}{The model and the material}
  \fullplate[0.94]{two_scales_question.pdf}
\end{frame}

% What the method is handed and what it is allowed to suppose, before any of it is built. Every
% line is the manuscript's own: the goal is the abstract's closing sentence
% (paper/sn-article.tex:637), what is fixed comes from "The sheet material, the outline, the
% topology and the actuation are fixed before the design begins" and what is free from "Two
% choices remain: where the cuts fall, and how wide the ligament left between them is", and the
% five assumptions are the \begin{hypotheses} list at paper/sn-article.tex:1621 in its own order.
%
% ⚠ THE ORDER OF THE FIVE IS THE MANUSCRIPT'S AND IS LOAD-BEARING. Coplanarity is second because
%   the Gamma criterion later on the deck is what discharges it, and the state-function assumption
%   is fourth because the annex's path-dependence figure is what measures its error. Renumbering
%   here silently renumbers what those two slides refer back to.
% ⚠ No numbers on this slide, deliberately. It is the one frame that states the problem rather
%   than a result, and a figure quoted here would have to be defended twice.
% ⚠ The three design freedoms are named here SEPARATELY and are not collected into theta. The
%   assembled design vector is drawn once, on the hinge frame, and stating it here as well would
%   be the same definition on two slides.
\begin{frame}{The problem}
  \punch{Given a target \symtgt{\mathcal{T}} and an actuation \symtgt{\zeta}, find where to cut
         the sheet so that it opens onto that target, while limiting hinge damage and the
         compression and shear that drive it.}
  \vspace{3mm}
  \begin{columns}[t,onlytextwidth]
    \begin{column}{0.40\textwidth}
      {\footnotesize\color{chartergrey}fixed before the design}
      \begin{itemize}\setlength{\itemsep}{1.2mm}
        \item the sheet material
        \item the outline and the topology
        \item the actuation, \symtgt{\zeta}
      \end{itemize}
      \vspace{3mm}
      {\footnotesize\color{chartergrey}left to the design}
      \begin{itemize}\setlength{\itemsep}{1.2mm}
        \item where the cuts fall, \symcut{r} and \symbnd{\tau}
        \item how wide the ligaments are, \symhng{w_{\mathrm{lig},h}}
      \end{itemize}
    \end{column}\hfill
    \begin{column}{0.56\textwidth}
      {\footnotesize\color{chartergrey}assumed}
      \begin{enumerate}\setlength{\itemsep}{1.2mm}
        \item The panels are rigid. All deformation localises in the hinges,
              \symhng{\bm\varepsilon_h}.
        \item The panels remain coplanar throughout deployment.
        \item Hinge energy is local:
              \symhng{\mathcal{W}=\textstyle\sum_h\mathcal{W}(\bm\varepsilon_h;\bm g_h)}.
        \item \symhng{\mathcal{W}} and damage are single-valued in
              \symhng{\bm\varepsilon_h}, not in the path taken to it.
        \item Deployment is quasi-static.
      \end{enumerate}
    \end{column}
  \end{columns}
\end{frame}

% ═══════════════════════════════════════════════════════════════════════════════════════════
\section{A differentiable design chain}
% ═══════════════════════════════════════════════════════════════════════════════════════════
\sectionslide{}

% The manuscript's Figure 1, recomposed for a screen: the cards in one row rather than two, built
% by the plate's own generator under --slide so there is one source of truth for the drawings.
%
% FOUR cards here, not five. This is the first sight of the pipeline and nothing in it has been
% explained, so the hinge energy has no card yet and the deployment is drawn as what it is at this
% point in the talk, a mechanism opening: no energy field on the hinges and no applied moment.
% --kinematic is what collapses the two. The full five-card row, with the energy field and its
% colour key, returns on the objective frame once every stage has been through.
% ⚠ The row always fills 408 bp, so the dropped card is absorbed by the other four rather than
%   leaving a hole: each is 81.8 bp wide here against 63.4 on the five-card row, and the two
%   plates are therefore NOT at one scale. They are never side by side, so that is the right
%   trade, but do not crop one and show it beside the other. The page arrangement cannot be shown here at all -- it is 437 x 425 bp against a
% 410 x 188 pt body, so it prints at 47 per cent and takes its 7 pt type to 3.3.
% The plate is 408 x 128 bp and \textwidth is 409.7 pt, so at [0.62] the WIDTH binds and it prints
% at 1.004 of the size it was authored at.
\begin{frame}{The pipeline}
  \fullplate[0.94]{pipeline_row_kin.pdf}
\end{frame}

% The deck showed WHAT is computed and never stated it. Every symbol here is the manuscript's
% own: Pi and the inner problem from eq:equilibrium, the outer problem from eq:objective, the
% implicit function theorem from eq:ift. The Introduction's own compact display pairs the two
% argmins the same way.
%
% \epsh is a paper macro (\bm{\varepsilon}) and does not exist here, so it is written out. There is
% no amssymb in this deck. Display skips are set locally and length is paid for in words: a
% negative \vspace big enough to buy a line pulls the next paragraph on top of the one above it
% and TeX reports ZERO overfull boxes.
\begin{frame}{From flat sheet to deployed shape}
  \fullplate[0.90]{shield_pipeline.pdf}
\end{frame}

% What the parameterisation gives, next to what sets it. Left: card (a) of the manuscript's
% Figure 1 in TWO crops, the sheet with its r and tau knobs above the zoom of one cut with rL and
% (1-r)L. Right: the same kind of sheet cut and then opened, PURELY KINEMATICALLY.
%
% The design-vector band that sits between those two crops in the source is dropped here and shown
% on the hinge frame instead, where g_h is defined and theta = (r, tau, {g_h}) can be read.
%
% Each crop is cut to its OWN content, not to a shared window, and the two are sized separately.
% There is no scale to hold between them: the lower box is labelled "zoom" in the source and is a
% magnified detail of one cut, so the sheet is free to be the larger of the two, which is the order
% they should be read in. A shared window would instead hold the sheet down to the width the wider
% zoom box needs, and leave the card's cream ground standing either side of it as a cut band.
% Content boxes, off a 300 dpi rasterisation of the 437.31 x 424.57 bp plate with 1.5 bp of air:
% sheet 76.44 x 76.20 bp (the orange square, its edge tabs and the tau label's overhang), zoom
% 104.04 x 49.32 bp (its own rounded border). At 113 pt and 93 pt the sheet prints 112.6 pt tall
% against the 69.9 pt it reached inside the shared window.
% ⚠ Rows with no pixel below 200 give the breaks. A lighter threshold catches the card's cream
%   ground and its container border and runs all three bands together.
%
% The kinematic pair is the 'chart' rung of the deployability harvest built on its own
% (--arms chart --one-scale). That rung is the closed chart with no equilibrium anywhere in it, so
% its flat sheet and its deployed sheet are a mechanism opening and nothing more.
%
% ⚠ THE TWO CELLS ARE AT ONE SCALE AND MUST STAY THERE. --one-scale draws them on one window, so
%   the sheet growing as it opens is a change of size and not two drawings each fitted to its own
%   cell: the ink measures 53.04 bp square flat and 85.68 bp square deployed, a linear expansion
%   of 1.615. Each is cropped to its own ink and then shown at 1.70 pt per bp, so 57.04 bp of crop
%   becomes 96.97 pt and 89.68 becomes 152.46. Equal WIDTHS or equal HEIGHTS here would both throw
%   the scale away; what has to be equal is pt per bp.
% Boxes measured by scanning a 300 dpi rasterisation of the 130.08 x 288.96 bp build with 2 bp of
% air added around each. Both come out square to three decimals, which is the check that the crop
% is on the sheet and not on something beside it.
% ⚠ SKIP THE LEFT GUTTER AT x < 25 bp WHEN MEASURING. The rotated row labels "flat sheet" and
%   "deployed" sit there. Included, they carry the ink box out to x = 12 and the trims that follow
%   cut the left edge off both sheets.
% ⚠ RE-MEASURE THESE FOUR TRIMS WHENEVER kinematic_pair IS REBUILT. The generator's left margin
%   moves with its readout layout, which slides both sheets, and a stale trim shows as a sheet with
%   a straight edge sliced off it rather than as any kind of build error.
%
% \mbox around the row: it cannot break. The row came to 431.6 pt against a 432.4 pt measure once
% and TeX silently wrapped it onto a second line, which surfaced as a 117 pt overfull vbox and said
% nothing about the cause. Unbreakable, an overrun names the width instead.
\begin{frame}{The cut pattern}
  \bled{\mbox{%
    \begin{minipage}[b]{113pt}\centering
      \includegraphics[trim={22.00bp 316.53bp 338.87bp 31.84bp},clip,%
                       width=\linewidth]{pipeline_overview.pdf}\\[1.6mm]
      \includegraphics[trim={8.08bp 236.37bp 325.19bp 138.88bp},clip,%
                       width=93pt]{pipeline_overview.pdf}
    \end{minipage}%
    \hspace{20pt}%
    \begin{minipage}[b]{96.97pt}\centering
      \includegraphics[trim=45.32bp 208.48bp 27.72bp 23.44bp,clip,%
                       width=\linewidth]{kinematic_pair.pdf}\\[1.4mm]
      {\scriptsize\color{chartergrey}flat sheet}
    \end{minipage}%
    \hspace{13pt}\raisebox{50pt}{\color{chartergrey}\large$\rightarrow$}\hspace{13pt}%
    \begin{minipage}[b]{152.46pt}\centering
      \includegraphics[trim=29.00bp 84.64bp 11.40bp 114.64bp,clip,%
                       width=\linewidth]{kinematic_pair.pdf}\\[1.4mm]
      {\scriptsize\color{chartergrey}deployed}
    \end{minipage}%
  }}
  \punch{$r$ void aspect ratio.\quad $\tau$ boundary position.\quad
         $\bm g_h$ hinge design vector.}
\end{frame}

% The author's own photograph. Rotation is baked into the pixels and the EXIF tag stripped:
% \includegraphics does not read EXIF, and straight off the phone this lands a quarter turn
% over. See assets/SOURCES.md for the conversion and the crop.
% This is where the talk turns from geometry to mechanics. A rigid-panel pattern assumes its panels
% stay flat and coplanar. Cut in PAPER at this panel size they do not: the panels themselves buckle
% out of plane, which is the coplanarity assumption failing rather than the hinges tearing.
% ⚠ NOT card, and NOT a hinge failure. Corrected on the author's instruction: the sheet is paper and
%   what is visible is out-of-plane buckling of the panels. This is the same phenomenon the Gamma
%   ladder measures deliberately, Gamma = (L/L*)^3 below one, where the paper puts the panels
%   tilting out of plane with their ligaments (sec:buckling-locality).
% 1312 x 678 px, aspect 1.94, against a body of 2.29, so the height binds.
% No band and no title: the picture and one line, and nothing else. It is meant to land as a pause
% before the mechanics starts. [plain] drops the headline and the footline; giving no frametitle is
% what stops the band, since the band is the frametitle template. The frame still advances the
% counter, so every number after it is unchanged.
\begin{frame}[plain]
  \photoplate[0.86]{whynottakingmaterialintoaccountfails.jpg}
  \punch{Cut in paper, the panels buckle out of plane.}
\end{frame}

% The design vector spelled out, in the manuscript's own order and words (its glossary entry at
% paper/sn-article.tex:921).
% ⚠ alpha carries (kappa) on purpose. The cut angle is NOT free: it follows from the cut pattern,
%   so of the five only the ligament width is set per hinge, and the fillet ratio, the kerf and the
%   window radius are shared. Dropping the argument would assert five independent knobs per hinge.
%
% The design-vector band beside the plate is the manuscript's own, cropped out of card (a) of
% Figure 1 rather than set again here, so there is one drawing of theta = (r, tau, {g_h}) and not
% two. It reads here and not on the parametrisation frame because g_h is what this frame defines:
% the braces are the whole set of hinge descriptors, and the punch line under it says what one
% descriptor holds.
% Geometry: the band's ink is 76.24 x 30.64 bp with 2 bp of air, shown 100 pt wide, so it prints
% at 1.31 of the size it was authored at. The plate is 370.6 x 188.12 bp, so at the 0.81 height cap
% the HEIGHT binds and it comes out 302 pt wide: 302 + 18 + 100 = 420 on the 432.4 pt bled measure.
% The raisebox centres the band on the plate rather than dropping it to the plate's baseline.
\begin{frame}{The hinge}
  \bled{\mbox{%
    \includegraphics[width=\charterplatew,height=0.81\charterbodyht,keepaspectratio]%
                    {hinge_layout_slide.pdf}%
    \hspace{18pt}%
    \raisebox{56.5pt}{\includegraphics[trim={21.98bp 284.59bp 339.09bp 109.34bp},clip,%
                                       width=100pt]{pipeline_overview.pdf}}%
  }}
  \punch{$\bm g_h = (w_{\mathrm{lig},h},\ \alpha_h(\bm\kappa),\ \rho_f,\ w_c,\ r_{\mathrm{win}})$:
         ligament width, cut angle, fillet ratio, kerf, window radius.}
\end{frame}

% The author's own photograph. Rotation is baked into the pixels and the EXIF tag stripped:
% \includegraphics does not read EXIF, and straight off the phone this lands a quarter turn
% over. See assets/SOURCES.md for the conversion and the crop.
% Two measurements, one per photograph: the coupons the material law is fitted to, and one ligament
% opened through a prescribed rotation with the buckle visible. Both portrait, so both are set to a
% common HEIGHT and centred, with a bare label under each.
% Only pictures, so no band and no title. The two captions say what each one is, which is all the
% frame needs. Heights are against \textheight for the same reason as the frame above.
\begin{frame}[plain]
  \bled{\mbox{%
    \begin{minipage}[b]{196pt}\centering
      \includegraphics[height=0.90\textheight,keepaspectratio]{meusingtheinstron.jpg}\\[1.4mm]
      {\scriptsize\color{chartergrey}tensile coupons, force and extension}
    \end{minipage}%
    \hspace{22pt}%
    \begin{minipage}[b]{196pt}\centering
      \includegraphics[height=0.90\textheight,keepaspectratio]{measuringoutofplanedeformation.jpg}\\[1.4mm]
      {\scriptsize\color{chartergrey}one ligament, through a prescribed rotation}
    \end{minipage}%
  }}
\end{frame}

% One PET ligament in pure relative rotation, across the range of widths solved. That reading
% is in the claim line rather than a subtitle: the plate is 370.6 x 220.3 bp against a body
% 398.3 x 196.3, so it is height-bound, and a subtitle costs it seven per cent of its size.
\begin{frame}{The learned hinge law}
  \fullplate[0.86]{energy_landscape.pdf}
  \punch{Past $40^\circ$: median deflection \textbf{53.8} sheet thicknesses.}
\end{frame}

\begin{frame}{Equilibrium}
  \begin{columns}[c,onlytextwidth]
    \begin{column}{0.27\textwidth}
      \centering
      \pipec[0.45]
    \end{column}\hfill
    \begin{column}{0.69\textwidth}
      \small
      % Displays, not \vspace: a negative skew big enough to buy a line pulls the NEXT paragraph
      % on top of the one above it, and TeX reports no overfull box because the space it was
      % given back makes the column "fit". Set the display skips and pay for length in words.
      \setlength{\abovedisplayskip}{1.5mm}\setlength{\belowdisplayskip}{1.5mm}%
      \setlength{\abovedisplayshortskip}{1.1mm}\setlength{\belowdisplayshortskip}{1.1mm}%
      At load factor $\zeta$ the sheet minimises
      \[
        \Pi(\bm q;\bm\theta,\zeta)=
        \underbrace{\sum_{h}\mathcal{W}\bigl(\bm\varepsilon_h(\bm q);\bm g_h\bigr)}
                   _{\text{stored in the hinges}}
        \;-\;\zeta\,\underbrace{\sum_{k}\bigl(\bm F_k\!\cdot\!\bm u_k+M_k\phi_k\bigr)}
                   _{\text{work of the applied loads}}
      \]
      Inner problem, the deployed state:
      \[ \bm q^\star(\zeta;\bm\theta)=\arg\min_{\bm q\in\mathcal{C}}\ \Pi(\bm q;\bm\theta,\zeta) \]
      Outer problem, the design:
      \[ \bm\theta^\star=\arg\min_{\bm\theta\in\Theta}\ \mathcal{L}(\bm\theta),
         \qquad \mathcal{L}=\mathcal{L}_{\mathrm{geom}}+\mathcal{L}_{\mathrm{phys}} \]
      Design gradient, from the converged state:
      \[ \partial\bm q^\star\!/\partial\bm\theta
         =-H^{-1}\,\partial^2\Pi/\partial\bm q\,\partial\bm\theta,
         \qquad H=\nabla^2_{\bm q}\Pi \]
    \end{column}
  \end{columns}
\end{frame}

% The pipeline reduced, with what the loop is actually minimising written under it. It closes the
% methods part: the plate is the one the deck opened the section with, and by here every stage of
% it has been explained.
%
% \epsh is a paper macro and does not exist in this deck, so the hinge deformation measure is
% written out as \bm\varepsilon_h, which is what the equilibrium frame uses. Every symbol and both
% sums are the manuscript's own, from eq:objective at paper/sn-article.tex:2478.
% ⚠ The asymmetry between the two terms is the point of the slide and is stated: the shape term is
%   evaluated on the LAST state alone, the physical term over the WHOLE path.
\begin{frame}{The objective}
  \fullplate[0.46]{pipeline_row.pdf}
  \setlength{\abovedisplayskip}{1.2mm}\setlength{\belowdisplayskip}{1.2mm}%
  \setlength{\abovedisplayshortskip}{0.9mm}\setlength{\belowdisplayshortskip}{0.9mm}%
  {\small
   \[ \mathcal{L}(\bm\theta)
      = \mathcal{L}_{\mathrm{geom}}\bigl(\mathcal{P}(\bm q^\star_N);\,\mathcal{T}\bigr)
      + \mathcal{L}_{\mathrm{phys}}\bigl[\{\bm\varepsilon_h(\bm q^\star_i)\}_{i,h}\bigr] \]}
  \begin{columns}[t,onlytextwidth]
    \begin{column}{0.55\textwidth}
      {\footnotesize
       $\mathcal{L}_{\mathrm{geom}}$, bidirectional Chamfer, on the deployed outline alone
       \[ \tfrac{1}{|\mathcal{P}|}\textstyle\sum_{\bm p}\min_{\bm t}\lVert\bm p-\bm t\rVert^{2}
          \;+\; w_{\mathrm{cov}}\,\tfrac{1}{|\mathcal{T}|}
          \textstyle\sum_{\bm t}\min_{\bm p}\lVert\bm t-\bm p\rVert^{2} \]
       \vspace{-1mm}
       precision, then coverage}
    \end{column}\hfill
    \begin{column}{0.41\textwidth}
      {\footnotesize
       $\mathcal{L}_{\mathrm{phys}}=\sum_j w_j\Phi_j$, over the whole path
       \[ \Phi_{\mathrm{damage}},\quad \Phi_{\mathrm{comp}},\quad \Phi_{\mathrm{shear}} \]
       \vspace{-1mm}
       already returned by the incremental solve}
    \end{column}
  \end{columns}
\end{frame}

% ═══════════════════════════════════════════════════════════════════════════════════════════
\section{Results and validation}
% ═══════════════════════════════════════════════════════════════════════════════════════════
\sectionslide{}
\subsection{Numerical experiments}

% The five problems side by side, one column each, with the three visuals that set them up: the
% loading diagram, the optimised cut pattern flat, and the deployed sheet against its target. The
% force transmission is NOT here. Each problem gets its own frame for that.
%
% Transposed in LaTeX rather than authored transposed in the generator. Every cell is a trim/clip
% out of the SAME matrix_all.pdf, so all fifteen sit at one scale and one type size and there is no
% second composition to keep in step with the first. Cells print at about 0.78 of authored size,
% which thins the strokes: that is the safe direction for a plate meant to be scanned before the
% five frames that read it closely.
%
% Geometry from the generator's own layout: gutter 12, panel_w 109.94, schem_w 68.17, head 6,
% row pitch 79.61, on the 432 x 456.04 bp canvas.
% ⚠ These are the --no-headers --no-row-labels numbers and transfer to no other build.
\begin{frame}{Five problems, one sheet}
  \bled{%
    \begin{tabular}{@{}*{5}{c@{\hspace{7pt}}}@{}}
      \includegraphics[trim=12.00bp 370.43bp 351.83bp 6.00bp,clip,height=0.242\charterbodyht]{matrix_all.pdf} &
      \includegraphics[trim=12.00bp 290.82bp 351.83bp 85.61bp,clip,height=0.242\charterbodyht]{matrix_all.pdf} &
      \includegraphics[trim=12.00bp 211.21bp 351.83bp 165.21bp,clip,height=0.242\charterbodyht]{matrix_all.pdf} &
      \includegraphics[trim=12.00bp 131.61bp 351.83bp 244.82bp,clip,height=0.242\charterbodyht]{matrix_all.pdf} &
      \includegraphics[trim=12.00bp 52.00bp 351.83bp 324.43bp,clip,height=0.242\charterbodyht]{matrix_all.pdf} \\[1.2mm]
      \includegraphics[trim=80.17bp 370.43bp 241.89bp 6.00bp,clip,height=0.242\charterbodyht]{matrix_all.pdf} &
      \includegraphics[trim=80.17bp 290.82bp 241.89bp 85.61bp,clip,height=0.242\charterbodyht]{matrix_all.pdf} &
      \includegraphics[trim=80.17bp 211.21bp 241.89bp 165.21bp,clip,height=0.242\charterbodyht]{matrix_all.pdf} &
      \includegraphics[trim=80.17bp 131.61bp 241.89bp 244.82bp,clip,height=0.242\charterbodyht]{matrix_all.pdf} &
      \includegraphics[trim=80.17bp 52.00bp 241.89bp 324.43bp,clip,height=0.242\charterbodyht]{matrix_all.pdf} \\[1.2mm]
      \includegraphics[trim=190.11bp 370.43bp 131.94bp 6.00bp,clip,height=0.242\charterbodyht]{matrix_all.pdf} &
      \includegraphics[trim=190.11bp 290.82bp 131.94bp 85.61bp,clip,height=0.242\charterbodyht]{matrix_all.pdf} &
      \includegraphics[trim=190.11bp 211.21bp 131.94bp 165.21bp,clip,height=0.242\charterbodyht]{matrix_all.pdf} &
      \includegraphics[trim=190.11bp 131.61bp 131.94bp 244.82bp,clip,height=0.242\charterbodyht]{matrix_all.pdf} &
      \includegraphics[trim=190.11bp 52.00bp 131.94bp 324.43bp,clip,height=0.242\charterbodyht]{matrix_all.pdf} \\[1.2mm]
      {\scriptsize\color{chartergrey}I} & {\scriptsize\color{chartergrey}II} & {\scriptsize\color{chartergrey}III} & {\scriptsize\color{chartergrey}IV} & {\scriptsize\color{chartergrey}V}
    \end{tabular}}
  \punch{150 panels, 275 hinges. Identical supports and target throughout.}
\end{frame}

\begin{frame}{I. Concentrated axial load}
  \designstrip{1}
  \bled{\mdrow{6.0bp}{763.94bp}}
  \punch{\forcekey}
\end{frame}

\begin{frame}{II. Transverse load, far corner}
  \designstrip{2}
  \bled{\mdrow{183.5bp}{586.26bp}}
  \punch{\forcekey}
\end{frame}

\begin{frame}{III. Couple on the loaded panel}
  \designstrip{3}
  \bled{\mdrow{361.18bp}{410.58bp}}
  \punch{\forcekey}
\end{frame}

\begin{frame}{IV. Eccentric axial load}
  \designstrip{4}
  \bled{\mdrow{536.86bp}{230.65bp}}
  \punch{\forcekey}
\end{frame}

\begin{frame}{V. Self-equilibrated shear couple}
  \designstrip{5}
  \bled{\mdrow{716.79bp}{52.67bp}}
  \punch{\forcekey}
\end{frame}

% ═══════════════════════════════════════════════════════════════════════════════════════════
\subsection{Design and validation on polyester sheet}
% ═══════════════════════════════════════════════════════════════════════════════════════════

% The author's own photograph. Rotation is baked into the pixels and the EXIF tag stripped:
% \includegraphics does not read EXIF. See assets/SOURCES.md.
% The turn from simulated actuations to a sheet on the floor. Left: a ligament in the polyester
% itself. Right: two-panel specimens folded by hand, each annotated with its own hinge width, panel
% length and Gamma, which is that quantity measured across scales.
% Gamma is what decides whether the panels stay coplanar, so this is also why the sheet about to be
% tested can be modelled with rigid panels at all.
%
% eq:gamma at paper/sn-article.tex:3738 is a moment balance about the hinge axis of one panel:
% gravity resisting the tilt against the folded ligament driving it. The RATIO is what the slide
% states, and nothing more.
% ⚠ Neither closed form goes on the slide. The expression 6 p L^3 w_lig/(E t^4) with p = rho g t
%   was tried and cut for being unreadable at a glance, and the reduced form (L/L*)^3 needs L*
%   defined before it says anything. A ratio of two moments needs neither and is what the quantity
%   IS. The algebra is in the paper for anyone who asks.
% ⚠ THE NUMBER IS CONDITIONAL AND THE CONDITIONS ARE PART OF IT. The paper's 3x10^2 is for PET at
%   w_lig = 18 mm AND a 406.4 mm pitch. Quoting it against the pitch alone leaves out the ligament
%   width it also depends on.
% ⚠ Gamma > 1 is the whole point: below it the panels tilt out of plane with their ligaments, which
%   is what the paper photograph two sections earlier shows happening. Without that the number is
%   just a magnitude.
\begin{frame}{The material, by hand}
  \bled{\mbox{%
    \begin{minipage}[b]{128pt}\centering
      \includegraphics[height=0.74\charterbodyht,keepaspectratio]{arealhinge.jpg}\\[1.4mm]
      {\scriptsize\color{chartergrey}one ligament, in the sheet}
    \end{minipage}%
    \hspace{24pt}%
    \begin{minipage}[b]{123pt}\centering
      \includegraphics[height=0.74\charterbodyht,keepaspectratio]{calibratingthehingedomain.jpg}\\[1.4mm]
      {\scriptsize\color{chartergrey}$\Gamma$ by hand, across scales}
    \end{minipage}%
  }}
  \punch{$\Gamma = \mathcal{M}_{\mathrm{grav}}/\mathcal{M}_{\mathrm{el}}$, and the panels stay
         coplanar for $\Gamma>1$.\quad Polyester, $w_{\mathrm{lig}}=18\,\mathrm{mm}$ at a
         $406.4\,\mathrm{mm}$ pitch: $\Gamma\approx3\times10^{2}$.}
\end{frame}

% The two designs one after the other, at one scale, each on the full body width. Clicked rather
% than set side by side: the comparison is between two rows of the same plate, and half a slide
% each leaves both at 30 per cent of the size they were authored at. See \hpdrowone above for the
% measured cut and for why row 1 carries no colour bar.
\begin{frame}{The design solve}
  \only<1>{\hpdrowone}
  \only<2>{%
    \hpdrowtwo
    \punch{Inside the training domain $60.2 \to \textbf{95.9\%}$.\quad
           Damage $0.175 \to \textbf{0.033}\,D_s$.\quad
           Ligament width $17.4\pm2.4 \to \textbf{21.6}\pm0.7$\,mm.}}
\end{frame}

\begin{frame}{From cut file to deployed sheet}
  \begin{columns}[c,onlytextwidth]
    \begin{column}{0.34\textwidth}
      \centering
      \includegraphics[width=\linewidth,height=0.94\charterbodyht,keepaspectratio]{sheet_closed_photo.png}
    \end{column}\hfill
    \begin{column}{0.62\textwidth}
      \setlength{\tabcolsep}{0pt}%
      \begin{tabular}{@{}p{34mm}l@{}}
        stock            & 4\,ft $\times$ 8\,ft polyester, $t = 0.5$\,mm\\[3mm]
        pattern          & $3\times3$, cut by hand\\[3mm]
        ligaments        & twelve, at the optimised widths\\[3mm]
        measurement      & photogrammetry, three states\\[3mm]
        floor            & 7.6\,mm on centres, $0.3^\circ$ on opening
      \end{tabular}
    \end{column}
  \end{columns}
\end{frame}

% The plate here is the top row of the merged validation plate: the deployed sheet in plan, and the
% one ringed hinge photographed alone against its prediction. It carries the sheet as well as the
% ligament, so the sheet-level agreement belongs on it. The state-by-state comparison those two
% numbers come from is in the annex.
\begin{frame}{Measured on the deployed sheet}
  \physhinge[0.75]
  \punch{Out of plane \textbf{23.7\,mm}, 47 sheet thicknesses.}
\end{frame}

% The agreement itself, moved out of the backup frames and into the talk: the chain is only worth
% anything if the sheet it predicts is the sheet that was built, and that claim needs its figure
% rather than a number quoted off one.
%
% This band and the hinge band on the frame above are two crops of ONE plate, read top to bottom,
% so the two frames sit at one scale. The plate's third band, the per-state charts, is a backup
% frame: the shapes ARE the agreement and carry it visually, and the charts quantify it for
% anyone who asks.
% ⚠ 17.1 and 0.70 are on THIS claim line and nowhere else in the talk. They were on the hinge
%   frame's line until the shapes came out of the annex, and putting them back there would state
%   both of them twice.
\begin{frame}{Shape agreement, state by state}
  \bled{\physstates}
  \punch{Panel corners \textbf{17.1 $\pm$ 3.6\,mm}, \textbf{0.70\%} of the sheet length.}
\end{frame}

% The photograph is the author's own; rotation baked into the pixels, EXIF tag stripped, since
% \includegraphics does not read EXIF. See assets/SOURCES.md.
%
% An expanded-metal grate is a steel sheet slit and stretched into a lattice, so it is the paper's
% own closing outlook already built: a deployable element cut flat from one sheet. What stands
% between that and this work is whether such a sheet's panels stay coplanar, which is Gamma.
% ⚠ Gamma is DEFINED on the frame that introduces it, "The material, by hand". It is not restated
%   here. The 0.79 is the mild-steel value from paper/sn-article.tex:3771, and it is the number
%   that says a steel sheet does NOT satisfy the assumption the polyester one does.
% ⚠ At 1.5 mm. The thickness is part of the number: Gamma goes as t^-4 through the elastic moment
%   and t^1 through the areal weight, so quoting 0.79 for "mild steel" without it is meaningless.
\begin{frame}{Conclusion}
  \begin{columns}[c,onlytextwidth]
    \begin{column}{0.60\textwidth}
      \begin{itemize}\setlength{\itemsep}{2.8mm}
        \item One differentiable chain: a material law from tensile coupons, a condensed hinge
              energy carrying buckling, yielding and damage, the equilibrium of the assembled sheet.
        \item Valid at both scales by construction. One network evaluation per hinge, about eight
              orders of magnitude faster than the solve it stands in for.
        \item Matching through equilibrium halves the boundary error. Penalising the hinge
              deformations raises the load carried by \textbf{64\%}.
        \item Cut from polyester and opened by hand: \textbf{0.70\%} of the sheet length.
      \end{itemize}
    \end{column}\hfill
    \begin{column}{0.36\textwidth}
      \centering
      \includegraphics[width=\linewidth,height=0.805\charterbodyht,keepaspectratio]%
                      {realworldkirigami.jpg}\\[1.4mm]
      {\scriptsize\color{chartergrey}expanded metal.\\
       mild steel at $1.5\,\mathrm{mm}$: $\Gamma = 0.79$}
    \end{column}
  \end{columns}
\end{frame}

{\setbeamertemplate{footline}{}
\begin{frame}[plain,noframenumbering]
  \charterfield
  \vfill
  \begin{center}
    {\color{charterink}\fontsize{30}{34}\selectfont\bfseries Thank you\par}
    \vskip4mm
    {\color{charterink}\large\itshape Questions\par}
  \end{center}
  \vfill
\end{frame}}

% ═══════════════════════════════════════════════════════════════════════════════════════════
% Backup frames, held for questions. They sit after the closing slide, so nothing here is
% presented unless it is asked for, and appendixnumberbeamer restarts the frame number at 1 so
% the count printed on a presented frame is the count of presented frames.
%
% Ordered as the questions come: why deployability is built into the parameterisation instead of
% penalised, how closely the hinge openings themselves agree, why the geometric problem is
% underconstrained, that the design gradient is verified, and that the hinge law the whole chain
% rests on is checked against the finite-element solutions it was fitted to.
% ═══════════════════════════════════════════════════════════════════════════════════════════
\appendix

% Four implementations of one task at equal budget, 49 panels deployed to a circle. The setup is
% in the claim line rather than a subtitle, for the height: the plate is 370.6 x 274.5 bp on a
% 398.3 x 196.3 body.
%
% Held for questions rather than presented. The talk asserts at the cut pattern that the
% parameterisation makes every pattern deployable by construction; this is the evidence, and it
% is the answer to "why not just penalise the invalid ones" rather than a result of its own. It is
% the first backup frame because that question arises earliest.
% ⚠ DO NOT crop the numbers row off the bottom to buy more. The two group boxes are anchored at
%   gs[2,0].y0 + 0.022 in plot_validity_ladder.py, so they enclose all three rows, and a cut
%   through the numbers leaves both boxes open-ended. That is the same reason
%   plot_physical_validation.py is run with --no-boxes before its rows are cut.
\begin{frame}{Four ways to keep a pattern deployable}
  \fullplate[0.96]{validity_ladder_slide.pdf}
\end{frame}

% The last band of the physical-validation plate, and the quantitative half of the agreement the
% talk makes visually two frames earlier. It carries its own title because it is no longer an
% overlay of the shapes frame, and it holds the RMSE number alone.
% ⚠ 0.3 degrees is the repeatability FLOOR of the measurement, not a model error. The claim is
%   that 2.4 is close to the floor, so dropping the floor would leave 2.4 sounding like a
%   shortfall instead of a limit.
\begin{frame}{Hinge opening, state by state}
  \bled{\physcharts}
  \punch{Hinge opening \textbf{2.4$^\circ$} RMSE, against a $0.3^\circ$ floor.}
\end{frame}

\begin{frame}{Two hundred optimisations of one target}
  \fullplate[0.875]{solution_set_slide.pdf}
  \punch{Two families, $41.5^\circ$ and $54.7^\circ$, \textbf{9.40}$\times$ apart.}
\end{frame}

\begin{frame}{Gradient checks}
  \charterplateheight
  \begin{center}
    \includegraphics[height=0.80\charterbodyht,keepaspectratio]{path_dependence_slide.pdf}\hspace{7mm}%
    \includegraphics[height=0.80\charterbodyht,keepaspectratio]{taylor_convergence_slide.pdf}
  \end{center}
  \punch{State function: \textbf{0.001\%} at the median hinge. Taylor remainder: order
         \textbf{2.05} constructed, \textbf{1.00} penalised.}
\end{frame}

\begin{frame}{Surrogate against finite elements}
  \fullplate[0.88]{surrogate_validity_slide.pdf}
  \punch{\textbf{1.70\%} worst, 0.97\% median, per state.}
\end{frame}

\end{document}
